% !TeX root = ../cs_thesis_main.tex

\chapter{Conclusions}

This study set out to address the critical operational gap in municipal solid waste management: the disconnect between static, ``one-size-fits-all'' policy frameworks and the non-linear behavioral realities of resource-constrained Local Government Units (LGUs). By departing from traditional administrative heuristics and developing a \textit{Heuristic-Guided Deep Reinforcement Learning (HuDRL)} framework coupled with a Multi-Level Agent-Based Model (ABM), this research successfully simulated the complex polycentric governance dynamics between municipal budgeting, social norms, and enforcement intensity in the Municipality of Bacolod, Lanao del Norte.

The simulation results provided a compelling computational argument to the ``dilution problem'' inherent in local governance. The study concluded that the administrative standard of ``Equal Distribution'' (the Status Quo) is \textit{structurally insufficient to overcome behavioral inertia in high-density zones}, validating the \textit{Urban Resource Trap} hypothesis where density and anonymity dissolve the social pressure that sustains compliance in rural communities \citep{Villanueva2021}. The HuDRL agent's independent discovery of the \textit{Sequential Saturation} strategy—which prioritizes the concentrated, asymmetric application of resources to trigger localized behavioral tipping points—offered a data-driven refutation of the status quo. It demonstrated that effective governance requires the political will to prioritize sequential efficacy over simultaneous, yet ineffectual, fairness.

\newpage

\section{Achieved Objectives}
\setlength{\parindent}{2em}

The study successfully met its primary research objectives, yielding critical exploratory insights into the socio-technical nature of waste segregation. First, the research successfully constructed a bounded, abstracted Agent-Based Model (ABM) parameterized with institutional data, such as municipal budgets and population distributions, alongside synthesized behavioral parameters derived from literature and key informant interviews. This computational environment is explicitly framed as an operational testbed for policy optimization, not a full ecological simulation. 

Second, regarding behavioral sensitivity within the model, the Global Sensitivity Analysis suggests that, under this specific parameterization, the Cost of Effort exerts the highest influence on simulated compliance variance ($S_1 = 0.79$). It is crucial to emphasize that this is a model-internal finding demonstrating within the computational environment that structural friction is the primary inhibitor of waste segregation, rather than a definitive claim about real-world causality \citep{Yazawa2025}. Furthermore, the Attitude parameter exhibited near-zero direct sensitivity, computationally quantifying the Value-Action Gap and confirming that passive Information, Education, and Communication (IEC) campaigns suffer severe diminishing returns when physical barriers remain unaddressed \citep{Zhao2022}.

Third, in terms of policy optimization within the model, the HuDRL agent successfully learned a dynamic allocation policy that consistently outperforms all static baselines across multiple metrics. The comparative analysis demonstrated within the model that traditional single-instrument regimes are structurally flawed: a ``Status Quo'' strategy acts as a dilution mechanism spreading funds too thinly, while ``Pure Enforcement'' and ``Pure Incentives'' rapidly trigger political collapse or budget exhaustion. Instead, the primary operational insight derived from the simulation is the agent's discovered ``Sequential Saturation'' strategy, which concentrates capital to trigger a social norm shield in targeted areas before reallocating resources.

Finally, providing a comparative evaluation, the optimized HuDRL policy achieved a terminal compliance of 90.45\% in the simulation. This stands in stark contrast to the 30.18\% simulated compliance achieved by the Status Quo approach, as well as the early systemic collapse triggered by the Pure Enforcement heuristic. Ultimately, these results are strictly exploratory and comparative; they demonstrate the methodology's potential as a robust decision-support tool for local government units, rather than providing definitive, guaranteed real-world prescriptions.

\section{Policy Brief}
\label{sec:policy_brief}

This section translates the computational findings from the Agent-Based Model and Deep Reinforcement Learning (ABM-DRL) simulation into an actionable policy brief for the Municipality of Bacolod, Lanao del Norte. The proposed Data-Driven Solid Waste Management (DD-SWM) Framework provides a strategic roadmap for municipal ordinance drafting, budget allocation, and enforcement design.

\subsection*{Executive Summary}
The Municipality of Bacolod faces a persistent challenge in enforcing household waste segregation under R.A. 9003. Despite sustained awareness campaigns, compliance remains critically low at approximately 12.5\%. Computational simulation of the municipality’s socio-technical dynamics reveals that the primary cause of failure is not a lack of total funding, but a structural flaw in how resources are distributed. The prevailing administrative heuristic of “Equal Distribution” dilutes municipal capital, preventing high-density areas from receiving the concentrated intervention required to trigger behavioral change.

By leveraging Artificial Intelligence to simulate thousands of policy permutations, this study identifies a mathematically optimal alternative: \textbf{Sequential Saturation}. The LGU must abandon simultaneous, equal budget distribution in favor of a phased rollout—concentrating heavy infrastructure and enforcement in one pilot barangay to force a localized “social tipping point,” then rapidly scaling back to a maintenance budget to rescue the next target.

\subsection*{The Challenge: Structural Flaws in Current SWM Policy}
Municipal audits and simulation baselines indicate an objective segregation compliance rate of approximately 12.5\%. The simulation identified two structural flaws driving this systemic non-compliance:

\begin{enumerate}
    \item \textbf{The Convenience Barrier (Value-Action Gap):} The simulation’s Global Sensitivity Analysis proved that the physical “Cost of Effort” is the absolute dominant bottleneck to compliance. Citizens are not failing to segregate due to a lack of environmental awareness; they are failing because the physical hassle of compliance remains insurmountably high. Passive Information, Education, and Communication (IEC) campaigns cannot overcome this logistical friction.
    \item \textbf{The Urban Resource Trap:} Allocating the \textpeso 1.5 million annual augmentation fund equally across all seven barangays ensures failure in dense urban centers like Barangay Poblacion. The budget per capita drops below the critical threshold necessary to maintain enforcement visibility and logistical support, leading to widespread behavioral apathy.
\end{enumerate}

\subsection*{Evidence from AI Simulation}
To test potential solutions, an AI agent (Heuristic-Guided Deep Reinforcement Learning, HuDRL) was tasked with managing the municipality’s budget over a simulated three-year horizon. The results were unambiguous:

\begin{itemize}
    \item \textbf{Status Quo (Equal Distribution):} Stagnated at roughly 30.18\% global compliance. Rural areas succeeded, but urban centers flatlined near 0\%.
    \item \textbf{Pure Enforcement:} Achieved temporary compliance spikes but rapidly collapsed due to severe political backlash and perceived unfairness.
    \item \textbf{Pure Incentives:} Created a “Transactional Trap” where funds were quickly exhausted, leading to a collapse in intrinsic motivation and compliance.
    \item \textbf{HuDRL (Sequential Saturation):} Reached over 90\% municipal-wide compliance without exceeding the \textpeso 375,000 quarterly budget by dynamically pulsing resources where they were most needed.
\end{itemize}

\subsection*{Actionable Policy Recommendations}
To operationalize the AI’s findings, the Sangguniang Bayan and the MENRO are advised to adopt the four pillars of the DD-SWM Framework.

\subsubsection*{\textit{Recommendation 1: Legislate a Phased Rollout (Sequential Saturation)}}
The Sangguniang Bayan must formally pass an appropriation ordinance that legally ring-fences municipal SWM funds for a sequential rollout. Instead of dividing the budget seven ways, 50–70\% of the quarterly budget should be heavily concentrated on a single prioritized “Pilot Zone” (e.g., Poblacion) for a fixed duration. After a barangay crosses 70\% compliance, its funding must be drastically titrated to a fractional “Maintenance Budget” (5–10\% of the quarterly fund) to free capital for the next target. To mitigate political accusations of bias, the LGU must use transparent, algorithmically derived criteria for prioritization and establish a publicly visible rotation schedule, with independent monitoring by civil society organizations.

\subsubsection*{\textit{Recommendation 2: Subsidize Friction Reduction over Passive IEC}}
The LGU must shift funding away from purely educational campaigns and toward direct subsidies that lower the physical effort of segregation. This includes:
\begin{itemize}
    \item Distributing standardized, color-coded receptacles to low-income households (Appendix B.6).
    \item Prioritizing the rapid repair of collection vehicles (e.g., the non‑functional MultiCab in Liangan East).
    \item Institutionalizing a purok-level collection system where localized waste workers gather waste directly from households while serving as frontline compliance monitors. Building on the successful grassroots practice in Barangay Binuni (Appendix B.3), collectors should be appointed as plantilla personnel with a regular salary, equipped with motorized cabs, and rotated quarterly to prevent bias.
\end{itemize}

\subsubsection*{\textit{Recommendation 3: Implement Pulsed Enforcement and the “Graduation Rule”}}
The MENRO must deploy its Municipal Strike Force and temporary incentives as a \textit{catalytic} tool, not a permanent fixture. During a barangay’s initial saturation phase, enforcement and incentives should be maximized. However, once compliance exceeds 70\%, the LGU must immediately invoke the “Graduation Rule”: scale back active enforcement and incentives to a fractional maintenance level. This prevents the “Cobra Effect” (where citizens only comply for a reward) and allows localized social norms to sustain the behavior. Any incentive program must be explicitly time‑limited with a clear phase‑out schedule.

\subsubsection*{\textit{Recommendation 4: Institutionalize Predictive Policy Prototyping}}
The LGU should not revert to drafting policies based purely on historical intuition. The municipality should partner with the Department of Science and Technology (DOST) to pilot a localized “Digital Twin” decision-support system based on the DRL architecture developed in this study. This tool will allow local executives to input demographic shifts, vehicle logistics, and budget constraints, mathematically stress‑testing SWM ordinances in a virtual environment before committing taxpayer funds.

\subsection*{Expected Outcomes and Next Steps}
Adopting the DD-SWM Framework shifts the Municipality of Bacolod from reactive enforcement to proactive, algorithmic governance. By legally adopting the Sequential Saturation strategy, the LGU can realistically achieve sustained compliance above 85\% across all jurisdictions within a three-year horizon, without requiring external budget loans.

\textbf{Immediate Next Steps:}
\begin{enumerate}
    \item The Municipal Solid Waste Management Board must convene to formally invalidate current equal-distribution models.
    \item The MENRO must select the first high-density barangay (e.g., Poblacion) for the Phase 1 Saturation pilot.
    \item The Sangguniang Bayan must draft an ordinance codifying the Rotation Schedule, the Graduation Rule, and the purok-level collector-inspector system.
    \item A memorandum of agreement with DOST should be initiated to develop the predictive policy prototyping tool.
\end{enumerate}

\subsection*{Acknowledgement of Political Feasibility}
The Sequential Saturation strategy is mathematically optimal within the simulation, but its real-world feasibility is constrained by the political economy of Philippine local governance, where patronage, pork barrel logics, and political dynasties often override algorithmic efficiency \citep{Carpio2025a,IBON2026,PIDS2024}. Therefore, the DD-SWM Framework should be interpreted as a \textit{proof of concept} that highlights the gap between computational optimality and on‑the‑ground political reality. Successful implementation would require prior institutional reforms—civil society oversight, anti‑dynasty legislation, and strengthened local accountability mechanisms—that lie beyond the scope of this thesis. Nevertheless, the framework provides a rigorous, evidence‑based starting point for pilot testing in a single barangay, from which political and operational lessons can be learned.


\section{Critique and Limitations}
\setlength{\parindent}{2em}

While the coupled ABM-DRL framework provides robust strategic insights, the model is subject to specific constraints that define the scope of its predictive applicability.

A primary limitation is the \textit{spatial abstraction} of the simulated environment. The current model utilizes a discrete grid and a network-based Moore neighborhood topology to represent barangays, rather than continuous Geographic Information System (GIS) coordinates. Consequently, the calculation of the \textit{Cost of Effort} assumes a uniform distance to collection points within a zone, averaging out micro-spatial factors such as road quality, elevation, or precise infrastructure placement that heavily influence real-world convenience.

Crucially, a fundamental limitation is inherent to the \textit{operational optimization framing} of the study. By design, this research prioritizes internal comparative validity over external predictive accuracy. The ABM is deliberately abstract, the behavioral parameters are synthesized from literature and qualitative interviews rather than empirically measured, and the multi-objective reward function encodes specific researcher choices. Consequently, the model cannot be used to forecast real-world compliance rates, specific waste diversion tonnages, or guaranteed cost-benefit ratios for the actual municipality. Its intrinsic value lies in demonstrating a computational methodology and generating data-driven hypotheses—such as the efficacy of sequential saturation and the dominance of physical effort—that could be tested in small-scale municipal pilots. Readers and policymakers are strictly cautioned against treating the specific numerical outputs, such as the 90.45\% simulated compliance rate, as realistic empirical predictions.

Furthermore, the model assumes a \textit{homogeneity of agent typologies}. Although agents possess distinct socioeconomic weights and income profiles, they all belong to a single functional class: ``Residential Households''. The simulation does not currently account for commercial entities (e.g., public markets) or institutional generators (e.g., schools), which likely operate under entirely different utility functions and may require distinct enforcement strategies compared to residential units.

Additionally, the \textit{Social Norm Shield} was implemented as a heuristic threshold (programmatically locking $w_{sn}$ values to 0.90 upon reaching the 70\% tipping point) rather than a fully emergent cognitive property. While this accurately aligns with sociological observations of habit formation and the ratchet effect \citep{Corcoran2020}, it remains a simplified computational representation of the complex psychological process of norm internalization. 

Finally, the simulation does not explicitly model the \textit{informal economy}. As qualitative interviews suggested, practices such as localized bribery or ``tipping'' collection personnel to bypass the ``no segregation, no collection'' rule can artificially lower the perceived cost of non-compliance, effectively functioning as an informal reduction in $c_{\text{effort}}$ and undermining the utility calculus established by the DRL strategy.

\section{Future Work}

To further refine the computational governance framework and transition it into a comprehensive municipal Digital Twin, future research should focus on the integration of \textit{GIS and Route Optimization}. By coupling the HuDRL agent with a Vehicle Routing Problem (VRP) solver using real-world shapefiles of Bacolod, the model could simultaneously optimize dynamic budget allocation and the physical routing of garbage trucks, providing a highly granular calculation of logistical costs \citep{Torkayesh2021}. Incorporating continuous spatial coordinates and vehicle capacity constraints will exponentially increase the state space dimensionality, necessitating advanced state-representation techniques such as Graph Neural Networks (GNNs) or multi-agent decentralized training paradigms to maintain computational tractability.

Future iterations should also introduce \textit{heterogeneous agent architectures}. Expanding the population to include distinct behavioral classes for commercial and industrial actors would allow the DRL agent to learn differential, sector-specific enforcement strategies—for example, optimizing strict corporate penalties while simultaneously subsidizing residential households.

Finally, the study is constrained by \textit{Validation Scope and Data Scarcity}. Due to the absence of longitudinal empirical datasets for out-of-sample testing, the framework relies on Face Validity and Structural Verification rather than Predictive Validity. While the HuDRL framework demonstrates internal consistency and theoretical alignment, its predictive accuracy regarding specific real-world tonnages remains untested. Future research should involve a pilot implementation in a single barangay to collect 'ground truth' data, which can then be used to perform Bayesian Melding or other empirical calibration techniques to move the model toward predictive validity.
\section{Final Remarks}

This thesis demonstrates that the challenge of Municipal Solid Waste Management is not merely a logistical engineering problem, but a complex, non-linear socio-technical dilemma. The HuDRL agent did not ``solve'' waste management through an influx of infinite resources; it solved it by revealing that the standard administrative intuition for ``fairness'' is mathematically inefficient in a resource-constrained system.

By illustrating that the binding constraint on R.A. 9003 compliance is not the total amount of money, but the spatial and temporal strategy by which it is deployed, this study offers the Municipality of Bacolod a scientifically grounded path forward. It signifies a necessary paradigm shift from reactive, intuition-based administration to proactive, predictive evidence-based policy design. The proposed Data-Driven Solid Waste Management (DD-SWM) Framework ensures that finite municipal resources are strategically invested where they can successfully overcome structural friction, triggering the self-sustaining social norms essential for long-term environmental resilience.

\noindent{\textbf{Usage of Generative AI}}

During the preparation of this work, the authors utilized Large Language Models (specifically Google Gemini) exclusively for the purpose of linguistic refinement and grammatical error correction. The generative AI tool was employed solely to enhance the clarity, coherence, and readability of the manuscript. No scientific content, data generation, experimental design, or conceptual analysis was performed by the AI; all intellectual contributions and final conclusions remain the sole work of the authors.